\scanpage{18}\mspage{17}
\deleted{\unsure{14|4}} \uline{15 juin}

On peut remplacer $\unsure{\mathcal F|\mathscr F}$ par divers\deleted{\unsure{es}} \inserted{\unsure{passage|pass. ge} à une} sous-ensembles ordonnés, stables \uline{par} \deleted{$\unsure{\leq|\in}$} \deleted{et} \deleted{par Sup}
sous-figures \deleted{fermées}\inserted{(i.e. \unsure{fermés|fermé} dans $\mathcal F$)} et par $\mathrm{Sup}$.

Exemples : Si $\unsure{L|\mathcal L}$ est un esp. top. \inserted{(str. top.)}, \deleted{de}
on peut soumettre les $A\in F$ \deleted{\unsure{figur}} ($F$ figure)
à des conditions de nature topologiques, comme d'être fermés, ou \deleted{\illegible{1}} compacts, ou fermés et connexes. Toute condition
portant sur la nature des strates $A\in F$ (en tant que partie de $L$, et non en tant que multistrate\ldots), ou $=$ sur la nature des
multistrates $F_A$ $(A\in F)$, définit une sous ens. ordonné de $\mathcal F$ ayant les propriétés de stabilité voulues. Signalons ici
des propriétés \unsure{portant} sur \unsure{la|le} $A^\circ=A\setminus\partial A$ (la définition de $A^\circ$ dépend de la multistrate $F_A$), comme celle d'être \deleted{non vide} \deleted{\unsure{connexé|connexe}},
\deleted{\unsure{ou}} \deleted{\unsure{ou}} connexe, ou plus gén. $i$-connexe $\unsure{0|-1}\leq i\leq+\infty$, \deleted{\unsure{$(-1)$-connexe signifiant non vide}} \inserted{et idem pour les $A$ eux-mêmes}. Ou on peut demander aux $|\partial A|$
d'être des sphères topologiques etc.
% slanted
\margin{\uline{NB} \deleted{Mais \unsure{l'axiome|l'énoncé} de \unsure{fi.|pr.} \unsure{supposé} d'\unsure{emblée}, -- peut-être \unsure{a-t-on|stat.}, ou \unsure{même} \unsure{ceci}, \unsure{ne|de} \unsure{saurait}-il \unsure{même} \unsure{convenir|voir} \unsure{en} \unsure{spécifiant} \unsure{cette} \unsure{notion} \unsure{avec|sur} des \unsure{fig.} \unsure{ou|un} $\unsure{F}$ \unsure{seulement} \uline{\unsure{préadmissibles}}.} Aussi, la condition de \uline{densité} $A=\overline{A^\circ}$ i.e. $A\subset B\Longrightarrow A^\circ\subset\overline{B^\circ}$.}

Si $L$ est \inserted{un espace} modéré, on se borne aux $\unsure{\mathscr{L}|\Phi}\subset\unsure{\mathcal P|\mathfrak P}(L)$ pour lesquels $A\in\mathscr{L}\Longrightarrow A$ partie fermée \uline{\unsure{par ordinaire|par ordre|pour ordre}} de $L$

On peut envisager aussi des conditions de nature plus globale sur $\mathscr{L}$, au cas où
